\lstset{
  mathescape=true,
  emph={[1]funcion, return, ret, retorno},
  emph={[2]para, cada, si, sino, if, else, true, false, mientras, not, o, y, null, en, de},
  emph={[3]agrego, borrar, copiar, agregar,sort,concatenar},
  emph={[4]from, to, head, next},
  emphstyle={[1]\color{violet}\it},
  emphstyle={[2]\color{red}\it},
  emphstyle={[3]\color{OliveGreen}\it},
  emphstyle={[4]\footnotesize\color{violet}},
  morecomment={[l]{//}},
  commentstyle={\color{gray}\it\footnotesize}
}

\def\grafo14{
  \begin{tikzpicture}[
      baseline = 0,
      grafo style,
      minimum size = 16pt,
      arista/.style={thick}
    ]
    \node[nodo] (v1) {$v_1$};
    \node[nodo, right of = v1] (v2) {$v_2$};
    \node[nodo, right of = v2] (v3) {$v_3$};
    \node[nodo, below right of = v1] (v4) {$v_4$};
    \node[nodo, right of = v4] (v5) {$v_5$};
    \node[nodo, below of = v4] (v6) {$v_6$};
    \node[nodo, right of = v3] (v7) {$v_7$};

    \draw[arista] (v1) to (v2) to (v3) to (v4);
    \draw[arista] (v1) to (v4) to (v5) to (v7);
    \draw[arista] (v4) to (v2) to (v5) to (v7);
    \draw[arista] (v6) to (v4);
    \draw[arista, bend right = 40pt] (v7) to (v2);
    \draw[arista, bend right = 40pt] (v6) to (v7);

    \node[above of = 1] {$G$};
  \end{tikzpicture}
}

\begin{enunciado}{\ejercicio[RepresentaGrafos*Hacer*]}
  \textit{RepresentaGrafos}

  Discutir (brevemente) las ventajas y desventajas (de las \textbf{\red{E}structuras de \red{D}atos}) en cuanto
  a la complejidad temporal y espacial de las siguientes implementaciones de un conjunto de vecindarios
  para un grafo $G$, de acuerdo a
  las siguientes operaciones:

  \textbf{Operaciones}
  \begin{enumerate}[label=$\alph*)$, itemsep=-2pt]
    \item Inicializar la estructura a partir de un conjunto de aristas de $G$.
    \item Determinar si dos vértices $v$ y $w$ son adyacentes.
    \item Recorrer y/o procesar el vecindario $N(v)$ de un vértice $v$ dado.
    \item Insertar un vértice $v$ con su conjunto de vecinos $N(v)$.
    \item Insertar una arista $vw$.
    \item Remover un vértice $v$ con todas sus adyacencias.
    \item Remover una arista $vw$.
    \item Mantener un orden de $N(v)$ de acuerdo a algún invariante que permita recorrer cada
          vecindario en un orden dado.
  \end{enumerate}

  \textbf{\red{E}structuras de \red{D}atos}
  \begin{enumerate}[label=$\text{\red{ED}}_\arabic*)$, itemsep=-2pt]
    \item $N$ se representa con una secuencia (vector o lista) que en cada posición $v$ tiene el
          conjunto $N(v)$ implementado sobre una secuencia (vector o lista). Cada vértice es una
          estructura que tiene un índice para acceder en $O(1)$ a $N(v)$. Esta representación se
          conoce comúnmente como \textit{lista de adyacencias}.

    \item Ídem anterior, pero cada $w \en N(v)$ se almacena junto con un índice a la posición que
          ocupa $v$ en $N(w)$. Esta representación también se conoce como \textit{lista de adyacencias},
          pero tiene información para implementar operaciones dinámicas.

    \item $N(v)$ se representa con un vector que en cada posición $i$ tiene un vector booleano $A_i$
          con $|V (G)|$ posiciones tal que $A_i[j]$ es verdadero si y solo si $i$ es adyacente a $j$. Esta
          representación se conoce comúnmente como \textit{matriz de adyacencias}.

    \item $N(v)$ se representa con un vector que en cada posición tiene el conjunto $N(v)$ implementado
          con una tabla de hash. Esta representación es un mix entre las representaciones clásicas de
          \textit{matriz de adyacencias} y \textit{lista de adyacencias}.
  \end{enumerate}
\end{enunciado}

\textit{Sobre las \red{E}structuras de \red{D}atos}:

Partiendo de un grafo \ul{no} dirigido ejemplificante $G$ voy a ir describiendo las distintas estructuras propuestas:
\begin{enumerate}[label=$\text{\red{ED}}_\arabic*)$, itemsep=-2pt]
  \item\label{ej-14:ed1} Un arreglo $A$ de tamaño $|V| = n$ para el cual el elemento $i$ apunta a una \textit{lista enlazada}
        cuyos nodos son el conjunto de vecinos del nodo $i$. Es decir que $A[i] \to N(v_i)$.

        \begin{minipage}{0.3\textwidth}
          $$
            \grafo14
          $$
        \end{minipage}
        \begin{minipage}{0.6\textwidth}
          $$
            \begin{tikzpicture}[
                scale=1,
                baseline = 0,
                rect/.style={ draw, thick},
                arrowstyle/.style={-latex, thick},
                submatriz/.style={ draw, very thick, anchor = center, inner sep = 2pt, color=#1, rounded corners=2pt,
                  },
              ]

              \matrix (A) [
                matrix of math nodes,
                nodes = { rect, column sep=-\pgflinewidth, minimum size = 23pt, },
              ]
              {
                1 & 2 & 3 & 4 & 5 & 6 & 7 \\
              };
              \foreach \x in {1,...,7}{
                  \node[above of = A-1-\x] () {$v_{\x}$};
                }

              \matrix (B) [matrix of math nodes,
                nodes={rect, rounded corners = 3pt},
                column sep=6pt, row sep=10pt,
                below=15pt of A,
              ]
              {
                v_2 & v_1 & v_2 & v_1 & v_2 & v_4 & v_2 \\
                v_4 & v_3 & v_4 & v_2 & v_4 & v_7 & v_5 \\
                    & v_4 &     & v_3 & v_7 &     & v_6 \\
                    & v_5 &     & v_5 &     &     &     \\
                    & v_7 &     & v_6 &     &     &     \\
              };

              \foreach \c in {1,...,7} {
                  \draw[arrowstyle] (B-1-\c) -- (B-2-\c);
                  \draw[arrowstyle] (A-1-\c) -- (B-1-\c);
                }

              \foreach \c in {2,4,5,7} {
                  \draw[arrowstyle] (B-2-\c) -- (B-3-\c);
                }

              \foreach \c in {2,4} {
                  \foreach \r \next in {3/4, 4/5}
                  \draw[arrowstyle] (B-\r-\c) -- (B-\next-\c);
                }

              \node[submatriz=BrickRed, dashed, fit=(A-1-1)(A-1-7), inner sep = 4pt, label={left, BrickRed}:Arreglo]{};

              \node[submatriz=Cerulean, fit=(B-1-1)(B-2-1), label={below, Cerulean}:\tiny$N(v_1)$](N1){};
              \node[submatriz=Cerulean, fit=(B-1-2)(B-5-2), label={below, Cerulean}:\tiny$N(v_2)$]{};
              \node[submatriz=Cerulean, fit=(B-1-3)(B-2-3), label={below, Cerulean}:\tiny$N(v_3)$]{};
              \node[submatriz=Cerulean, fit=(B-1-4)(B-5-4), label={below, Cerulean}:\tiny$N(v_4)$]{};
              \node[submatriz=Cerulean, fit=(B-1-5)(B-3-5), label={below, Cerulean}:\tiny$N(v_5)$]{};
              \node[submatriz=Cerulean, fit=(B-1-6)(B-2-6), label={below, Cerulean}:\tiny$N(v_6)$]{};
              \node[submatriz=Cerulean, fit=(B-1-7)(B-3-7), label={below, Cerulean}:\tiny$N(v_7)$]{};

              \node[label={left, Cerulean}:Listas Enlazadas] at (N1.west){};
            \end{tikzpicture}
          $$
        \end{minipage}
        \begin{enumerate}[label=$\alph*)$]
          \item Una lista de adyacencia que inicializo a partir de un conjunto de aristas tiene un \textit{running time} de:
                $$
                  \cajaResultado{
                    \en O(n+m)
                  }
                $$
                Dado que tengo que generar para cada uno de los $n$ vértices una lista pasando por las $m$
                aristas, creando el $N(v)$ en el proceso. Un pseudocódigo para eso:
                \begin{tcolorbox}
                  \begin{lstlisting}
para i en 1..n        //Inicializo el arreglo para cada vértice
    A[i] $\ot$ i

para arista en E[1..m]     // Ponele que arista = (from, to)
    A[arista.from] agrego arista.to
    A[arista.to] agrego arista.from \end{lstlisting}
                \end{tcolorbox}

                La complejidad espacial:
                $$
                  \cajaResultado{
                    \en O(n + m)
                  }
                $$
                Este último resultado para un grafo con $m\gg n \to O(n^2)$, complejidad espacial como la \textit{matriz de adyacencia}.

          \item Para saber si dos vértices $v$ y $u$ son adyacentes, es decir si existe una arista entre ellos tengo un \textit{running time}:
                $$
                  \en O(\ua{1}{\textit{\green{get}} \\ N} + \minimo(d(u), d(v)) \en O(V)
                $$
                En un grafo no dirigido la información de la arista está repetida, una vez en $A[u]$ y otra vez en $A[v]$
                entonces puedo buscarla \textit{en simultáneo} cortando cuando se termine la más corta:
                \begin{tcolorbox}
                  \begin{lstlisting}
vActual $\ot$ A[v].head
uActual $\ot$ A[u].head

mientras vActual.next y uActual.next not null // Corto cuando se
    si vActual=u o u.Actual=v              //terminó la vecindad más corta
        ret "son vecinos"
    sino 
        vActual $\ot$ vActual.next 
        uActual $\ot$ uActual.next 
ret "no son vecinos" \end{lstlisting}
                \end{tcolorbox}
                En el \textit{más peor} de los casos recorro todo el grafo.

          \item\label{ej-14:item-1-c}
                Nada muy loco por aquí, si $|N(v)| = d(v)$.
                $$
                  \cajaResultado{
                    \en O(1 + d(v))
                  }
                $$

          \item Agregar un vértice y sus vecinos, me obliga a llamar a todos los vecinos $w \en N(v)$ e insertar ahí a $v$, agrandar
                el arreglo de ser necesario también agregaría el costo de la nueva dimensión:
                $$
                  \cajaResultado{
                    \en \ub{O(1 + d(v))}{\text{con} \\ \text{espacio}} \en \ub{O(n + d(v))}{\text{sin}}
                  }
                $$
                \begin{tcolorbox}
                  \begin{lstlisting}
B[2*|A|]             // Declaro nuevo arreglo B con |B| > |A|. El doble en este ejemplo.
copiar(A,B)           // clono
A$\ot$B                // A tiene la misma info pero el doble de tamaño. La joda salió O(n)
para cada w en N(v) 
    agregar v en N(w)  // Update vecindades adyacentes a v \end{lstlisting}
                \end{tcolorbox}

          \item  Insertar una arista tendrá un costo constante. Si tengo que insertar una arista $vw$:
                $$
                  \cajaResultado{
                    (\textit{\green{get}} ~ N(v) ) + (\textit{\green{insertar w en}} ~ N(v)) \en O(1)
                  }
                $$

          \item Remover un vértice $w$ implica borrarlo \textit{de todos lados}, es decir tengo que buscarlo
                en cada elemento del arreglo.
                $$
                  \cajaResultado  {
                    \en O(n + m)
                  }
                $$
                \begin{tcolorbox}
                  \begin{lstlisting}
para cada v$\distinto$w en A[1..n]
    para cada vecino en N(v)      // O(d(v))
        si vecino=w 
            borrar vecino de N(v)
    // Actualizo A 
    A[v]$\ot$N(v)
    A[w] $\ot$ null\end{lstlisting}

                \end{tcolorbox}

          \item Remover una arista $vw$ es similar al anterior, pero ahora sé a cuales $N$ tengo que ir:
                $$
                  O(d(v) + d(w))
                $$
                \begin{tcolorbox}
                  \begin{lstlisting}
para cada vecino en N(e.from)      // e = (e.from, e.to) = (v,w) = v->w
    si vecino=e.to 
        borrar vecino de N(e.from)
e$\ot$(e.to, e.from) y repito   // la invierto y repito \end{lstlisting}
                \end{tcolorbox}

          \item \Hacer
        \end{enumerate}

        \bigskip

        \parrafoDestacado{
          \it
          Ahora en vez de usar listas enlazadas uso vectores, arreglos ordenados. Esta estructura es muy buena
          para hacer consultas en el grafo, dado el ordenamiento. Pero en contrapartida es
          una porquería para operaciones dinámicas, es decir, para modificar el grafo, agregando o sacando
          vértices, etc.
        }
        El \brick{arreglo} tiene en el elemento $i$ el valor del índice donde está primer vecino del vértice $i$ y consecutivo a $i$ los demás vecinos
        del vértice $i$, el cual tiene $d(v_i) = \brick{A}[i+1] - \brick{A}[i]$ vecinos.
        El orden en el \green{arreglo} de los vecinos me va a permitir buscar con menor complejidad que en la \textit{lista enlazada}.

        Por ejemplo $N(v_2) = \brick{A}[2] = 3 \to B[3..(3 + d(v_2) - 1)] = [1,3,4,5,7]$.

        Parecido a la \textit{lista de adyacencia}, esto es un \textit{arreglo de adyacencia}.

        \begin{minipage}{0.2\textwidth}
          \grafo14
        \end{minipage}
        \begin{minipage}{0.65\textwidth}
          $$
            \begin{tikzpicture}[
                baseline = 0,
                rect/.style={ draw, thick},
                arrowstyle/.style={-latex, thick},
                submatriz/.style={draw, anchor = center, thick, inner sep = 1pt, color=#1, rounded corners=4pt,
                  },
              ]

              \matrix (A) [
                matrix of math nodes,
                nodes = { rect, column sep=-\pgflinewidth, minimum size = 23pt, },
              ]
              {
                1 & 3 & 8 & 10 & 15 & 18 & 20 \\
              };
              \foreach \x in {1,...,7}{
                  \node[above of = A-1-\x] () {$v_{\x}$};
                }

              \matrix (B) [matrix of math nodes,
              nodes={rect, rounded corners = 0pt},
              column sep=1pt, row sep=10pt,
              below=15pt of A,
              header/.style={fill=lightgray!40},
              ]
              {
              |[header]| 2 &
              4 &
              |[header]|1 &
              3 &
              4 &
              5 &
              7 &
              |[header]|2 &
              4 &
              |[header]|1 &
              2 &
              3 &
              5 &
              6 &
              |[header]|2 &
              4 &
              7 &
              |[header]|4 &
              7 &
              |[header]|2 &
              5 &
              6 \\
              };

              \node[submatriz=BrickRed, dashed, fit=(A-1-1)(A-1-7), inner sep = 4pt, label={left, BrickRed}:Arreglo]{};
              \node[submatriz=orange,dashed, fit=(B-1-1)(B-1-22), inner sep = 4pt, label={left, orange}:Arreglo] {};
              \foreach \a/\b/\c in {1/1/2,2/3/5,3/8/2,4/10/5,5/15/3,6/18/2,7/20/3} {
                  \draw[arrowstyle] (A-1-\a.south) -- (B-1-\b.north);

                  \pgfmathsetmacro{\bmenosa}{\b+\c-1}
                  \node[submatriz=Cerulean, fit=(B-1-\b)(B-1-\bmenosa), label={below, Cerulean}:\tiny$N(v_{\a})$](){};
                }
            \end{tikzpicture}
          $$
        \end{minipage}

        \begin{enumerate}[label=$\alph*)$]
          \item Como el arreglo lo quiero ordenado, eso va a costar un $\log$ extra. La complejidad temporal queda en:
                $$
                  \cajaResultado{
                    \en O(n \cdot \log(n) + m)
                  }
                $$
                \begin{tcolorbox}
                  \begin{lstlisting}
A[n]        // Declaro arreglo A |V|
para arista en E[1..m]     // Ponele que arista = (from, to)
    ...                       //genero un arreglo N(from) = [to1, to2,...] para cada arista
para cada N(v)          // Ordeno esto que cuesta O(n*log(n))
    sort(N(v))
para i en [1..n] // Inicializo el vector A con valor A[v]=i tal que $A[i]=N_v[i]$
    si i = 1
        A[i]$\ot$1 
    sino
        A[i]$\ot$|N[i-1]|+1 \end{lstlisting}
                \end{tcolorbox}

          \item Acá es cuando el ordenamiento toma valor. El costo de saber si $v$ y $w$ son adyacentes será:
                $$
                  \cajaResultado{
                    \en O(\minimo(\log(d(u)), \log(d(v)))
                  }
                $$
                El grado de cada vértice lo saco en $O(1)$ del arreglo $A$. Por ejemplo:
                $$
                  d(i) = A[i+1] - A[i]
                $$

          \item  idem \ref{ej-14:ed1} con listas enlazadas.
          \item  Agregar un vértice $v$ me obliga a armar los arreglos $A$ y $B$ nuevamente, aunque me ahorraría
                el ordenamiento, ya que el nuevo vértice se pone al final del arreglo $A$ y al final de cada $N(v)$ de B.

                Luego restaría ordenar $N(v)$, lo que me agrega un $\log$. La complejidad quedaría:
                $$
                  \cajaResultado{
                    \en O(n + m + d(v) \cdot \log(d(v)))
                  }
                $$

          \item Agregar una arista es \textit{similar en costo} a agregar un vértice.
                \parrafoDestacado[\skull]{
                  Modificar esa estructura estaría siendo algo horrible.
                }
                Agregar una arista es agregar 2 vértices $v$ y $w$. Esto implica \textit{modificar, ordenando} dos
                grupos en $B$: $N(v) \ytext N(w)$, lo cual cuesta $O(d(v) \log(d(v)) + d(w) \log(d(w)))$
                Luego hay que recrear todo $B$ con los nuevos $N(v)$ y $N(w)$: $O(m)$. Un nuevo $B$ obliga
                a recalcular todo $A$: $O(n)$. El costo total queda:
                $$
                  \cajaResultado{
                    \en O(\ua{n}{A} + \ua{m}{B} +
                    \ub{d(v) \log(d(v))}{\text{agregar}\\\text{en} ~ N(v)}
                    +
                    \ub{d(w) \log(d(w))}{\text{agregar}\\\text{en} ~ N(w)})
                  }
                $$

          \item Remover un vértice con la \textit{lista de vectores} es parecido al \ref{ej-14:ed1} con listas enlazadas,
                porque tengo que construir todo $A$ y $B$ en peor caso, pero sin reordenar:
                $$
                  \cajaResultado{
                    \en O(n + m)
                  }
                $$

          \item Remover una arista $vw$ implica rearmar los vectores $A$ y $B$, dado que con el cambio del tamaño
                de $B$, cambian los valores de los elementos de $A$. Cómo no hay \textit{sorting} el costo es:
                $$
                  \cajaResultado{
                    \en O(n + m)
                  }
                $$

          \item \Hacer
        \end{enumerate}

  \item \Hacer

  \item \Hacer

  \item \Hacer
\end{enumerate}
